% Authors: Dr. Denys Dutykh (Mathematics Department, Khalifa University of Science and Technology,
% Abu Dhabi, UAE) and Prof. Laurent Vuillon (Univ. Savoie Mont Blanc, CNRS, LAMA, Chambery, France)
\section{Finite verifications and the proof of the theorem}\label{sec:assembly}

We now prove the main theorem. Its computer-assisted part consists of
finitely many assertions about integers and about elements of $\Z[\varphi]$;
we state them together, derive from them the three statements on which the
assembly rests, and then join the finite windows to the infinite tails at
explicit denominators, so that no interval between a finite calculation and
an asymptotic statement is left untreated. Throughout this section, for
$P=8,\ldots,13$ respectively, put
\begin{equation}\label{eq:tail-denominators}
 Q_P=29,35,39,76,101,146 .
\end{equation}
These are one beyond the six finite windows of \cref{tab:gaussian-windows}.
They satisfy $P+Q_P-3=c_P$, with $c_P$ as in
\eqref{eq:tail-source-windows}, and all six fractions $P/Q_P$ are reduced.
Beyond a second denominator the upper bound of \cref{prop:row-labels} alone
keeps a source below its ceiling $U_P$; with
\begin{equation}\label{eq:second-junctions}
 n_P=3,4,4,7,8,11,\qquad Q^*_P=P(n_P+1)=32,45,50,88,108,156,
\end{equation}
every reduced $P/Q$ with $Q\geqslant Q^*_P$ has runs of length at least $n_P$,
and the denominators $Q_P\leqslant Q<Q^*_P$ in between are finitely many.

\subsection{The finite verifications}

\begin{proposition}[Finite verifications]\label{prop:finite-verifications}
The following assertions hold.
\begin{enumerate}[label=\textup{(\roman*)},ref=\textup{(\roman*)},leftmargin=*]
\item\label{it:fv-windows} For each of the $248$ reduced fractions $P/Q$ of
 \cref{tab:gaussian-windows}, the label $m_{P/Q}$ is a product of explicit
 prime powers; each of these primes passes the criterion \eqref{eq:full-lucas}
 with explicit witnesses, recursively down to the prime two; each odd one is
 congruent to $1$ modulo $4$ and has an explicit decomposition as a sum of two
 squares; and among the pairs \eqref{eq:gaussian-sign-cube} built from these
 data, the occurrence pair of \cref{prop:occurrence-cost} has cost $P+Q-2$ and
 every other pair has a larger cost.
\item\label{it:fv-covers} For each $P\in\{8,\ldots,13\}$ and each $k$ with $2\leqslant k\leqslant P-6$,
 the subdivision tree of \cref{sec:tails} built from the root boxes of all
 $2^{k+1}$ parity sectors, with $S_0=c_P-k-1$ and the thresholds $C_P$ and
 $U_P$, is finite, and each of its leaves closes: its upper matched corner is
 strictly below $C_P$ or its lower matched corner is strictly above $U_P$.
 For $k=P-5$ and $S_0=c_P-k-1$, every root box has its upper matched corner
 strictly below $C_P$.
\item\label{it:fv-lower} For $k=2$ and $S_0=152$, every root box of the eight parity sectors has
 its upper matched corner strictly below $71/100$.
\item\label{it:fv-sources} For each $P\in\{8,\ldots,13\}$, one has
 $C_P\bigl(1+\varphi^{-4n_P-6}\bigr)^P<U_P$, and
 $m_{P/Q}<U_P\,F_{2(P+Q-2)+1}$ for each of the reduced $Q$ with
 $Q_P\leqslant Q<Q^*_P$, thirty-eight denominators in all.
\end{enumerate}
\end{proposition}

\begin{proof}
Each assertion is a finite statement about integers and about elements of
$\Z[\varphi]$, whose signs are decided exactly by the integer test of
\cref{sec:tails}. They were verified by computer in exact arithmetic, and the
verification was carried out twice independently, the two computations
resting on different formulas. For part~\ref{it:fv-windows}, the first computation builds the
fiber of each label from the sign cube \eqref{eq:gaussian-sign-cube} and
proves every prime by the criterion \eqref{eq:full-lucas}, applied
recursively down to the prime two; the second finds every representation
afresh, from the square roots of $-1$ modulo the label, combined by the
Chinese remainder theorem and reduced by Cornacchia's descent, and checks
every prime again, by the same criterion with independently written
arithmetic and by a probable-prime test. In both, each label is computed
from its Cohn matrix, the first computation obtaining it in two ways, from the
floor increments of the Christoffel word and by Stern--Brocot mediants. For
parts~\ref{it:fv-covers} and~\ref{it:fv-lower}, the first computation evaluates the normalized ratio
through the channel recursion \eqref{eq:tail-channel} in $\Z[\varphi]$; the
second rebuilds every tree from its root boxes and evaluates the ratio
through the spectral projectors of the matrix $\mathsf B$ in $\Z[\sqrt5]$. For
part~\ref{it:fv-sources}, the first computation takes each label from its Cohn matrix and
decides the constant comparison by the integer sign test in $\Q(\sqrt5)$;
the second takes each label as the continuant of the palindromic continued
fraction $[2,\sigma(z),2]$ of \c{C}anak\c{c}{\i} and Schiffler, the Fibonacci
numbers from the doubling formulas, and bounds the constants with rational
enclosures of $\sqrt5$ to sixty decimal places.
\end{proof}

For instance, the first integer comparison of part~\ref{it:fv-sources} is the positive
integer identity
$$
 707F_{71}-1000\,m_{8/29}=25054682255203.
$$

\begin{remark}\label{rem:computer-contribution}
The computer contributes only the assertions of
\cref{prop:finite-verifications}: the factorizations and the costs of
$1638$ explicit pairs, the closing of the boxes of finitely many trees,
thirty-eight comparisons of integers and six comparisons of constants.
Everything that turns them into statements about infinitely many
denominators is proved in the text: that the sign cube lists the whole
fiber (\cref{prop:gaussian-fibre}), that two matched corners bound a whole
box (\cref{thm:tail-parity} and \eqref{eq:tail-corner-bounds}), that the
root boxes cover every admissible gap tuple exactly once and each split
partitions a box (\cref{sec:tails}), that a source row stays in its band
beyond the junction (\cref{prop:row-labels} and \eqref{eq:second-junctions}),
and the comparisons at the remaining costs, below.
\end{remark}

\subsection{Three consequences}

\begin{proposition}[Complete finite windows]\label{prop:finite-windows}
For every reduced $P/Q$ in \cref{tab:gaussian-windows}, the occurrence
pair is the unique element of $\mathcal F(m_{P/Q})$ with minimum
Euclidean cost. That cost is $P+Q-2$.
\end{proposition}

\begin{proof}
By \cref{prop:finite-verifications}\ref{it:fv-windows} the label is a product of proved
primes, the odd ones split as recorded. By \cref{prop:gaussian-fibre}, the
pairs \eqref{eq:gaussian-sign-cube} built from this factorization are the
whole fiber $\mathcal F(m_{P/Q})$. By \cref{prop:finite-verifications}\ref{it:fv-windows}
again, the occurrence pair is the only one of them with cost $P+Q-2$, and
all others cost more. The table contains every reduced denominator of each
window, so no occurrence of the windows is missing.
\end{proof}

\begin{theorem}[Uniform tail exclusion]
\label{thm:tail-covers}
For each $P\in\{8,9,10,11,12,13\}$, each $2\leqslant k\leqslant P-6$, and
every nonnegative gap tuple with cost $c\geqslant c_P$, one has
$$
 \frac{N(g)}{F_{2c+3}}<C_P
 \quad\text{or}\quad
 \frac{N(g)}{F_{2c+3}}>U_P.
$$
The inequalities are strict, so no endpoint in these strata belongs to
the closed source band $[C_P,U_P]$. Every path with $k\geqslant P-5$
defects and cost $c\geqslant c_P$ satisfies the stronger conclusion
$N(g)/F_{2c+3}<C_P$.
\end{theorem}

\begin{proof}
\emph{Middle defect counts.} Fix $P$ and $k$ with $2\leqslant k\leqslant P-6$.
A tuple with $k$ defects has cost $c=S+k+1$, so $c\geqslant c_P$ means
$S\geqslant S_0=c_P-k-1$. By \cref{sec:tails}, the root boxes of all parity
sectors cover the tuples with $S\geqslant S_0$, each exactly once, and each
split partitions a box into two. By \cref{prop:finite-verifications}\ref{it:fv-covers} the resulting tree is
finite and every leaf closes, so each admissible tuple $g$ lies in exactly one
closed leaf. If the upper matched corner of that leaf is below $C_P$, then
\eqref{eq:tail-corner-bounds} gives
$\mathcal R_k(g)\leqslant\mathcal R_k(g^{\max})<C_P$; otherwise its lower
matched corner is above $U_P$ and
$\mathcal R_k(g)\geqslant\mathcal R_k(g^{\min})>U_P$.

\emph{The defect count $P-5$.} By \cref{prop:finite-verifications}\ref{it:fv-covers} every
root box has its upper matched corner below $C_P$, and
\eqref{eq:tail-corner-bounds} gives $\mathcal R_{P-5}(g)<C_P$ on the whole
root box, hence for every admissible tuple.

\emph{More defects.} If a path has more than $P-5$ defects, replace surplus
$\mathsf A$ moves by $\mathsf B$ as in \cref{lem:path-orders}. At an
ordered positive state $(u,v)$ the replacement increases the column by
$(v-u,0)^{\mathsf T}$, and every subsequent matrix is nonnegative and
invertible. The final norm strictly increases, while the cost stays fixed.
Replacing until exactly $P-5$ defects remain gives a path of the same cost
whose norm exceeds the original one and is below $C_PF_{2c+3}$.
\end{proof}

\begin{corollary}[A lower-cost bound]
\label{cor:tail-lower-cost}
Every path of cost $c\geqslant155$ with at least two defects satisfies
$$
 \frac{N(g)}{F_{2c+3}}<\frac{71}{100}.
$$
\end{corollary}

\begin{proof}
A path with exactly two defects and cost $c\geqslant155$ has $S\geqslant152$,
so it lies in one root box of the cover with $S_0=152$. By
\cref{prop:finite-verifications}\ref{it:fv-lower} and \eqref{eq:tail-corner-bounds}, its
normalized ratio is below the value at the upper matched corner of that box,
hence below $71/100$. A path with more defects is handled by replacing
surplus defects, as in the proof of \cref{thm:tail-covers}.
\end{proof}

The corollary is used only at the cost-$(L-2)$ interface in row thirteen.

\subsection{Where the source intervals meet the covers}

Put $L=P+Q-2$ and $c=L-1$. The occurrence pair has cost $L$, and we must
exclude every representation of the same norm at every smaller cost. We
first place the source in its band: for every reduced $P/Q$ with
$Q\geqslant Q_P$,
\begin{equation}\label{eq:joined-source-band}
 C_P<\frac{m_{P/Q}}{F_{2c+3}}<U_P.
\end{equation}
The lower bound is \cref{prop:row-labels}\ref{it:row-bounds}, since $F_{2c+3}=F_{2(P+Q-2)+1}$.
For the upper bound there are two cases. If $Q\geqslant Q^*_P$, every run of
$\gamma_{P/Q}$ has length at least $n=\lfloor(Q-P)/P\rfloor\geqslant n_P$, and since the
upper bound of \cref{prop:row-labels}\ref{it:row-bounds} decreases with $n$,
\begin{equation}\label{eq:junction-source}
 \frac{m_{P/Q}}{F_{2c+3}}<C_P\bigl(1+\varphi^{-4n_P-6}\bigr)^P<U_P ,
\end{equation}
the second inequality being part~\ref{it:fv-sources} of \cref{prop:finite-verifications}.
If $Q_P\leqslant Q<Q^*_P$, the upper bound is one of the thirty-eight integer
comparisons of the same part. \Cref{fig:junction} shows the band, the source
labels and the bound of \cref{prop:row-labels} for the row of numerator
eight.

\begin{figure}[tb]
\centering
% Authors: Dr. Denys Dutykh (Mathematics Department, Khalifa University of Science and
% Technology, Abu Dhabi, UAE) and Prof. Laurent Vuillon (Univ. Savoie Mont Blanc, CNRS, LAMA)
% Generated by supplement/code/make_figures.py from exact data; do not edit by hand.
\begin{tikzpicture}[font=\footnotesize, >={Stealth[length=4pt]}]
\fill[sand] (0,0) rectangle (12,3.649);
\draw[ochre, line width=0.8pt] (0,3.649) -- (12,3.649);
\node[ochre, above left, font=\scriptsize] at (12,3.649) {$U_8-C_8$};
\node[slate, below left, font=\scriptsize, align=right] at (11.95,3.599) {source band};
\draw[slate, densely dashed, line width=0.6pt] (3.243,0) -- (3.243,5.2);
\node[wine, anchor=south east, font=\scriptsize] at (3.193,5.2) {finite window};
\node[navy, anchor=south west, font=\scriptsize] at (3.293,5.2) {tail from $Q_8=29$};
\draw[teal, line width=0.7pt, dash pattern=on 5pt off 2pt] (0.081,4.972) -- (1.216,4.972) -- (1.216,4.474) -- (2.514,4.474) -- (2.514,4.011) -- (3.811,4.011) -- (3.811,3.552) -- (5.108,3.552) -- (5.108,3.095) -- (6.405,3.095) -- (6.405,2.637) -- (7.703,2.637) -- (7.703,2.18) -- (9,2.18) -- (9,1.722) -- (10.297,1.722) -- (10.297,1.264) -- (11.595,1.264) -- (11.595,0.807) -- (11.919,0.807);
\draw[teal, line width=0.5pt] (3.892,3.432) -- (3.892,3.672);
\node[teal, font=\scriptsize, inner sep=1pt] at (9.892,2.86) {bound of \cref*{prop:row-labels}};
\filldraw[wine, fill=white, line width=0.6pt] (0.162,4.731) circle (1.3pt);
\filldraw[wine, fill=white, line width=0.6pt] (0.486,4.658) circle (1.3pt);
\filldraw[wine, fill=white, line width=0.6pt] (0.811,4.561) circle (1.3pt);
\filldraw[wine, fill=white, line width=0.6pt] (1.135,4.411) circle (1.3pt);
\filldraw[wine, fill=white, line width=0.6pt] (1.459,4.251) circle (1.3pt);
\filldraw[wine, fill=white, line width=0.6pt] (1.784,4.185) circle (1.3pt);
\filldraw[wine, fill=white, line width=0.6pt] (2.108,4.095) circle (1.3pt);
\filldraw[wine, fill=white, line width=0.6pt] (2.432,3.949) circle (1.3pt);
\filldraw[wine, fill=white, line width=0.6pt] (2.757,3.791) circle (1.3pt);
\filldraw[wine, fill=white, line width=0.6pt] (3.081,3.726) circle (1.3pt);
\filldraw[navy, fill=white, line width=0.6pt] (3.405,3.636) circle (1.3pt);
\filldraw[navy, fill=white, line width=0.6pt] (3.73,3.49) circle (1.3pt);
\filldraw[navy, fill=white, line width=0.6pt] (4.054,3.333) circle (1.3pt);
\filldraw[navy, fill=white, line width=0.6pt] (4.378,3.268) circle (1.3pt);
\filldraw[navy, fill=white, line width=0.6pt] (4.703,3.178) circle (1.3pt);
\filldraw[navy, fill=white, line width=0.6pt] (5.027,3.033) circle (1.3pt);
\filldraw[navy, fill=white, line width=0.6pt] (5.351,2.875) circle (1.3pt);
\filldraw[navy, fill=white, line width=0.6pt] (5.676,2.81) circle (1.3pt);
\filldraw[navy, fill=white, line width=0.6pt] (6,2.721) circle (1.3pt);
\filldraw[navy, fill=white, line width=0.6pt] (6.324,2.575) circle (1.3pt);
\filldraw[navy, fill=white, line width=0.6pt] (6.649,2.418) circle (1.3pt);
\filldraw[navy, fill=white, line width=0.6pt] (6.973,2.353) circle (1.3pt);
\filldraw[navy, fill=white, line width=0.6pt] (7.297,2.263) circle (1.3pt);
\filldraw[navy, fill=white, line width=0.6pt] (7.622,2.117) circle (1.3pt);
\filldraw[navy, fill=white, line width=0.6pt] (7.946,1.96) circle (1.3pt);
\filldraw[navy, fill=white, line width=0.6pt] (8.27,1.895) circle (1.3pt);
\filldraw[navy, fill=white, line width=0.6pt] (8.595,1.805) circle (1.3pt);
\filldraw[navy, fill=white, line width=0.6pt] (8.919,1.66) circle (1.3pt);
\filldraw[navy, fill=white, line width=0.6pt] (9.243,1.502) circle (1.3pt);
\filldraw[navy, fill=white, line width=0.6pt] (9.568,1.437) circle (1.3pt);
\filldraw[navy, fill=white, line width=0.6pt] (9.892,1.348) circle (1.3pt);
\filldraw[navy, fill=white, line width=0.6pt] (10.216,1.202) circle (1.3pt);
\filldraw[navy, fill=white, line width=0.6pt] (10.541,1.045) circle (1.3pt);
\filldraw[navy, fill=white, line width=0.6pt] (10.865,0.98) circle (1.3pt);
\filldraw[navy, fill=white, line width=0.6pt] (11.189,0.89) circle (1.3pt);
\filldraw[navy, fill=white, line width=0.6pt] (11.514,0.745) circle (1.3pt);
\filldraw[navy, fill=white, line width=0.6pt] (11.838,0.587) circle (1.3pt);
\draw[ink, line width=0.5pt] (0,0) -- (12,0);
\draw[ink, line width=0.5pt] (0,0) -- (0,5.2);
\draw[ink, line width=0.5pt] (0.324,0) -- (0.324,-0.08);
\node[ink, below] at (0.324,-0.08) {$10$};
\draw[ink, line width=0.5pt] (1.946,0) -- (1.946,-0.08);
\node[ink, below] at (1.946,-0.08) {$20$};
\draw[ink, line width=0.5pt] (3.568,0) -- (3.568,-0.08);
\node[ink, below] at (3.568,-0.08) {$30$};
\draw[ink, line width=0.5pt] (5.189,0) -- (5.189,-0.08);
\node[ink, below] at (5.189,-0.08) {$40$};
\draw[ink, line width=0.5pt] (6.811,0) -- (6.811,-0.08);
\node[ink, below] at (6.811,-0.08) {$50$};
\draw[ink, line width=0.5pt] (8.432,0) -- (8.432,-0.08);
\node[ink, below] at (8.432,-0.08) {$60$};
\draw[ink, line width=0.5pt] (10.054,0) -- (10.054,-0.08);
\node[ink, below] at (10.054,-0.08) {$70$};
\draw[ink, line width=0.5pt] (11.676,0) -- (11.676,-0.08);
\node[ink, below] at (11.676,-0.08) {$80$};
\node[ink] at (6,-0.75) {denominator $Q$};
\draw[ink, line width=0.5pt] (0,0.274) -- (-0.08,0.274);
\node[ink, left] at (-0.08,0.274) {$10^{-9}$};
\draw[ink, line width=0.5pt] (0,1.368) -- (-0.08,1.368);
\node[ink, left] at (-0.08,1.368) {$10^{-7}$};
\draw[ink, line width=0.5pt] (0,2.463) -- (-0.08,2.463);
\node[ink, left] at (-0.08,2.463) {$10^{-5}$};
\draw[ink, line width=0.5pt] (0,3.558) -- (-0.08,3.558);
\node[ink, left] at (-0.08,3.558) {$10^{-3}$};
\draw[ink, line width=0.5pt] (0,4.653) -- (-0.08,4.653);
\node[ink, left] at (-0.08,4.653) {$10^{-1}$};
\node[ink, rotate=90] at (-1.25,2.6) {$m_{8/Q}/F_{2Q+13}-C_8$};
\fill[paper] (0.18,0.15) rectangle (3.05,3.45);
\draw[slate, line width=0.4pt] (0.18,0.15) rectangle (3.05,3.45);
\node[ink, font=\scriptsize, anchor=north west] at (0.22,3.43) {linear scale};
\draw[ochre, line width=0.6pt] (1.04,1.933) -- (2.87,1.933);
\node[ochre, above left, font=\tiny, inner sep=1.5pt] at (2.87,1.933) {$U_8-C_8$};
\draw[slate, densely dashed, line width=0.5pt] (1.65,0.48) -- (1.65,3.05);
\filldraw[wine, fill=white, line width=0.5pt] (1.345,2.918) circle (0.9pt);
\filldraw[navy, fill=white, line width=0.5pt] (1.955,1.789) circle (0.9pt);
\filldraw[navy, fill=white, line width=0.5pt] (2.565,0.661) circle (0.9pt);
\draw[ink, line width=0.4pt] (1.04,0.48) -- (2.87,0.48);
\draw[ink, line width=0.4pt] (1.04,0.48) -- (1.04,3.05);
\draw[ink, line width=0.4pt] (1.345,0.48) -- (1.345,0.42);
\node[ink, below, font=\tiny, inner sep=1.5pt] at (1.345,0.42) {$27$};
\draw[ink, line width=0.4pt] (1.955,0.48) -- (1.955,0.42);
\node[ink, below, font=\tiny, inner sep=1.5pt] at (1.955,0.42) {$29$};
\draw[ink, line width=0.4pt] (2.565,0.48) -- (2.565,0.42);
\node[ink, below, font=\tiny, inner sep=1.5pt] at (2.565,0.42) {$31$};
\draw[ink, line width=0.4pt] (1.04,1.1) -- (0.98,1.1);
\node[ink, left, font=\tiny, inner sep=1.5pt] at (0.98,1.1) {$0.001$};
\draw[ink, line width=0.4pt] (1.04,2.873) -- (0.98,2.873);
\node[ink, left, font=\tiny, inner sep=1.5pt] at (0.98,2.873) {$0.002$};
\end{tikzpicture}

\caption{The source band of the row $P=8$. The normalized labels
$m_{8/Q}/F_{2(Q+6)+1}$ (navy circles) decrease with the denominator and tend
to $C_8$; from the junction $Q_8=29$ on they lie in the band between $C_8$ and
$U_8=707/1000$ (shaded), which contains no normalized norm of a path one
step cheaper: \cref{thm:tail-covers} excludes the paths with at least two
defects, and the paths with fewer are excluded at the start of
\cref{sec:one-step-below}. Before the junction, the wine
circles are the labels of the finite window, settled by their complete
Gaussian fibers. The dashed staircase is the upper bound of
\cref{prop:row-labels}, $C_8\bigl((1+\varphi^{-4n-6})^8-1\bigr)$ with
$n=\lfloor(Q-8)/8\rfloor$; it lies above every label and enters the band at
$Q^*_8=32$, so that only the denominators $29$ and $31$ need a direct
comparison. The vertical scale is $(m_{8/Q}/F_{2(Q+6)+1}-C_8)$, on a
logarithmic axis. The inset repeats the denominators $27$, $29$ and $31$ on a
linear scale: at the junction the label lies inside the band by
$8.1\times10^{-5}$.}
\label{fig:junction}
\end{figure}

\subsection{Excluding cost one step below the occurrence}\label{sec:one-step-below}

The path with no defect has norm $F_{2c+3}$, above every source ceiling.
A one-defect path satisfies
$$
 \frac{N}{F_{2c+3}}\geqslant
 \frac{5F_{2c-1}}{F_{2c+3}}
 >5\varphi^{-4}>\frac{29}{40}>U_8\geqslant U_P.
$$
The first inequality is \cref{lem:one-defect-minimum}; the strict
Fibonacci comparison follows from Binet's formula at odd indices.
\Cref{thm:tail-covers} excludes the closed source band for every
remaining defect count. Thus no representation at cost $L-1$ has norm
$m_{P/Q}$.

\subsection{Excluding all smaller costs}

It remains necessary to check smaller costs: exclusion at one depth
alone would not imply minimality. We use the following elementary
consequences of Binet's formula. For positive odd $n$, the ratios
$F_n/F_{n+2}$ decrease strictly to $\varphi^{-2}$, and
$F_n/F_{n+4}$ decrease strictly to $\varphi^{-4}$. In particular,
for odd $n\geqslant5$,
\begin{equation}\label{eq:coarse-fibonacci-ratios}
 \varphi^{-2}<\frac{F_n}{F_{n+2}}\leqslant\frac5{13}<\frac25,
 \qquad
 \frac{F_n}{F_{n+4}}\leqslant\frac5{34}<\frac3{20}.
\end{equation}
All uses below occur strictly beyond these initial indices.

For $8\leqslant P\leqslant11$, every cost $d\leqslant L-2$ has
$$
 N\leqslant F_{2L-1}
 <\frac25F_{2L+1}
 <C_{11}F_{2L+1}\leqslant C_PF_{2L+1}\leqslant m_{P/Q}.
$$
For $P=12,13$, the same Fibonacci bound gives, for every cost $d\leqslant L-3$,
$$
 N\leqslant F_{2L-3}
 <\frac3{20}F_{2L+1}
 <C_{13}F_{2L+1}\leqslant m_{P/Q}.
$$
Only the cost $d=L-2$ in the last two rows remains. Its Fibonacci
maximum lies above the source, since
$$
 \frac{F_{2d+3}}{F_{2d+5}}>\varphi^{-2}>\frac38>U_{12}>U_{13}.
$$
For row twelve, every non-Fibonacci path satisfies
$N\leqslant E_d<F_{2d+5}/3<C_{12}F_{2d+5}\leqslant m_{12/Q}$.
Here the strict envelope comparison has a short exact proof:
$$
 F_{2d+5}-3E_d=2(F_{d+1}^2-2F_d^2)>0.
$$
Indeed $F_{d+1}/F_d>3/2$ for $d\geqslant5$, and the present range
starts at $d=109$.

For row thirteen, a one-defect path is above the source:
$$
 \frac{N}{F_{2d+5}}\geqslant\frac{5F_{2d-1}}{F_{2d+5}}
 >5\varphi^{-6}>\frac5{18}>U_{13}.
$$
The radical inequality uses $\varphi^6=9+4\sqrt5<18$.
Every path with at least two defects is below the source. Since
$d\geqslant155$, \cref{cor:tail-lower-cost} gives
$$
 \frac{N}{F_{2d+5}}
 =\frac{N}{F_{2d+3}}\frac{F_{2d+3}}{F_{2d+5}}
 <\frac{71}{100}\frac5{13}
 =\frac{71}{260}<C_{13}.
$$
All the comparisons made in this section between $C_P$, the rational
constants and $U_P$ are inequalities in $\Q(\sqrt5)$, decided exactly by
the integer sign test of \cref{sec:tails}; no decimal rounding enters the
proof.

\begin{proof}[Proof of \cref{thm:thirteen-rows}]
\Cref{prop:occurrence-cost} supplies a primitive representation
of cost $L=P+Q-2$. \Cref{prop:rows-seven} proves its weak
minimality for $P\leqslant7$. For each $8\leqslant P\leqslant13$,
\cref{prop:finite-windows} covers every reduced denominator
$P<Q<Q_P$. For $Q\geqslant Q_P$, the preceding three subsections
exclude every smaller cost. This proves weak minimality without a
denominator bound.
\end{proof}

\begin{proof}[Proof of \cref{cor:thirteen-injective}]
By \cref{thm:thirteen-rows}, every occurrence with numerator at most
thirteen is weakly Euclidean-minimal, and \cref{cor:restricted-injectivity}
applies to this family.
\end{proof}

The proof uses complete Gaussian fibers only for the finite initial
windows. In the infinite tails it excludes all shorter primitive paths
directly, so no factorization of an unbounded family of Markov integers
is required. These two forms of completeness are complementary.
